% GENERATED FILE. Edit canonical lessons, then run npm run generate:books.
% canonical-source-hash: fnv1a64:fa6977c9ddc1edac
% canonical-lessons: ZH-C445-hailang, ZH-C445-yuchuan, ZH-C445-shangan, ZH-C445-beike, ZH-C445-tuichao

\chapter{By the Sea}
\label{ch:zh-by-the-sea}
\hlchaptermodality{445}

\begin{chapteropening}
\textbf{By the end of this chapter:} \emph{I can talk about the waves, fishing boats, going ashore, shells and the tide going out.}

Complete the last lesson of chapter 445: I can talk about the waves, fishing boats, going ashore, shells and the tide going out.
\end{chapteropening}

\section[hǎilàng]{hǎilàng — sea wave; a wave of the sea}

\label{lesson:ZH-C445-hailang}



\begin{quote}
Before the new one: say the Mandarin for a lighthouse, then the Mandarin for sunscreen.
\end{quote}

\subsection*{You'll want to know: hǎilàng}
\emph{hǎilàng} — \textquotedblleft{}sea wave; a wave of the sea\textquotedblright{}.

\begin{quote}
\emph{hǎilàng tài dà le chuán kāi bu dào nàge dǎo} — The waves are too big: the boat can't get to the island.
\end{quote}

Its characters come later; for now it is a word for the ear.

\begin{cousinweb}
\emph{hǎilàng} is built from \zh{海} 'the sea' and \emph{làng} 'a wave', both inherited from Middle Chinese, \emph{xojX} and \emph{langH} (Guangyun, in Baxter's transcription), which glosses the second 'waves'.
\end{cousinweb}

\begin{grammarlens}[title={What you've built}]
One more word for getting around.
\end{grammarlens}

\subsection*{Your turn}
\begin{itemize}
\raggedright
  \item \emph{Say it:} \emph{hǎilàng}
  \item \emph{Say it:} \emph{hǎilàng}, once more
  \item \emph{Say it:} \emph{hǎilàng tài dà le chuán kāi bu dào nàge dǎo} — The waves are too big: the boat can't get to the island.
  \item \emph{Recall:} say the Mandarin for a lighthouse, then the Mandarin for sunscreen, then say \emph{hǎilàng} again
\end{itemize}

\begin{tcolorbox}[breakable,colback=teal!4,colframe=teal!35!black,title={Before you move on}]
What does hǎilàng mean? (\textquotedblleft{}Sea wave; a wave of the sea\textquotedblright{}.) Say it once more.
\end{tcolorbox}
\section[yúchuán]{yúchuán — fishing boat}

\label{lesson:ZH-C445-yuchuan}



\begin{quote}
Before the new one: say the Mandarin for sunscreen, then the Mandarin for a wave of the sea.
\end{quote}

\subsection*{You'll want to know: yúchuán}
\emph{yúchuán} — \textquotedblleft{}fishing boat\textquotedblright{}.

\begin{quote}
\emph{měitiān zǎoshang hěn duō yúchuán chūhǎi} — Every morning lots of fishing boats go out to sea.
\end{quote}

Its characters come later; for now it is a word for the ear.

\begin{cousinweb}
\emph{yúchuán} is built from \emph{yú} 'fishing' and \zh{船} 'a boat', both inherited from Middle Chinese, \emph{ngjo} and \emph{zywen} (Guangyun, in Baxter's transcription), which glosses the first, after an older dictionary, as catching fish. It is the syllable of \emph{yú} 'a fish', written with the water piece.
\end{cousinweb}

\begin{grammarlens}[title={What you've built}]
One more word for getting around.
\end{grammarlens}

\subsection*{Your turn}
\begin{itemize}
\raggedright
  \item \emph{Say it:} \emph{yúchuán}
  \item \emph{Say it:} \emph{yúchuán}, once more
  \item \emph{Say it:} \emph{měitiān zǎoshang hěn duō yúchuán chūhǎi} — Every morning lots of fishing boats go out to sea.
  \item \emph{Recall:} say the Mandarin for sunscreen, then the Mandarin for a wave of the sea, then say \emph{yúchuán} again
\end{itemize}

\begin{tcolorbox}[breakable,colback=teal!4,colframe=teal!35!black,title={Before you move on}]
What does yúchuán mean? (\textquotedblleft{}Fishing boat\textquotedblright{}.) Say it once more.
\end{tcolorbox}
\section[shàng'àn]{shàng'àn — to go ashore}

\label{lesson:ZH-C445-shangan}



\begin{quote}
Before the new one: say the Mandarin for a wave of the sea, then the Mandarin for a fishing boat.
\end{quote}

\subsection*{You'll want to know: shàng'àn}
\emph{shàng'àn} — \textquotedblleft{}to go ashore\textquotedblright{}.

\begin{quote}
\emph{xiàwǔ sì diǎn wǒmen shàng'àn zuò chē huí jiǔdiàn} — At four in the afternoon we go ashore and take the bus back to the hotel.
\end{quote}

Its characters come later; for now it is a word for the ear.

\begin{cousinweb}
\emph{shàng'àn} is built from \zh{上} 'to go up' and \emph{àn} 'a bank', both inherited from Middle Chinese, \emph{dzyangX} and \emph{nganH} (Guangyun, in Baxter's transcription), which glosses the first 'to climb' and the second as the high edge of the water.
\end{cousinweb}

\begin{grammarlens}[title={What you've built}]
One more word for getting around.
\end{grammarlens}

\subsection*{Your turn}
\begin{itemize}
\raggedright
  \item \emph{Say it:} \emph{shàng'àn}
  \item \emph{Say it:} \emph{shàng'àn}, once more
  \item \emph{Say it:} \emph{xiàwǔ sì diǎn wǒmen shàng'àn zuò chē huí jiǔdiàn} — At four in the afternoon we go ashore and take the bus back to the hotel.
  \item \emph{Recall:} say the Mandarin for a wave of the sea, then the Mandarin for a fishing boat, then say \emph{shàng'àn} again
\end{itemize}

\begin{tcolorbox}[breakable,colback=teal!4,colframe=teal!35!black,title={Before you move on}]
What does shàng'àn mean? (\textquotedblleft{}To go ashore\textquotedblright{}.) Say it once more.
\end{tcolorbox}
\section[bèiké]{bèiké — seashell}

\label{lesson:ZH-C445-beike}



\begin{quote}
Before the new one: say the Mandarin for a fishing boat, then the Mandarin for to go ashore.
\end{quote}

\subsection*{You'll want to know: bèiké}
\emph{bèiké} — \textquotedblleft{}seashell\textquotedblright{}.

\begin{quote}
\emph{wǒ zhǎodào le hěn duō bèiké fàng zài bāo lǐ} — I found lots of shells and put them in my bag.
\end{quote}

Its characters come later; for now it is a word for the ear.

\begin{cousinweb}
\emph{bèiké} is built from \emph{bèi} 'a shellfish' and \emph{ké} 'a hard shell'. The first is inherited from Middle Chinese \emph{pajH} (Guangyun, in Baxter's transcription), glossed as a creature of the sea in a shell; the second is in the Guangyun only under an older spelling, read \emph{khaewk} and glossed 'a hard skin', and the everyday \emph{ké} does not follow from that reading, so its history is not settled.
\end{cousinweb}

\begin{grammarlens}[title={What you've built}]
One more word for getting around.
\end{grammarlens}

\subsection*{Your turn}
\begin{itemize}
\raggedright
  \item \emph{Say it:} \emph{bèiké}
  \item \emph{Say it:} \emph{bèiké}, once more
  \item \emph{Say it:} \emph{wǒ zhǎodào le hěn duō bèiké fàng zài bāo lǐ} — I found lots of shells and put them in my bag.
  \item \emph{Recall:} say the Mandarin for a fishing boat, then the Mandarin for to go ashore, then say \emph{bèiké} again
\end{itemize}

\begin{tcolorbox}[breakable,colback=teal!4,colframe=teal!35!black,title={Before you move on}]
What does bèiké mean? (\textquotedblleft{}Seashell\textquotedblright{}.) Say it once more.
\end{tcolorbox}
\section[tuìcháo]{tuìcháo — to ebb (of the tide)}

\label{lesson:ZH-C445-tuichao}



\begin{quote}
Before the new one: say the Mandarin for to go ashore, then the Mandarin for a seashell.
\end{quote}

\subsection*{You'll want to know: tuìcháo}
\emph{tuìcháo} — \textquotedblleft{}to ebb (of the tide)\textquotedblright{}.

\begin{quote}
\emph{xiàwǔ tuìcháo de shíhou wǒmen qù hǎibiān zhǎo bèiké} — In the afternoon, when the tide goes out, we look for shells on the shore.
\end{quote}

Its characters come later; for now it is a word for the ear.

\begin{cousinweb}
\emph{tuìcháo} is built from \emph{tuì} 'to draw back' and \emph{cháo} 'the tide', both inherited from Middle Chinese, \emph{thwojH} and \emph{drjew} (Guangyun, in Baxter's transcription), which glosses the first 'to retreat' and the second as the tide.
\end{cousinweb}

\begin{grammarlens}[title={What you've built}]
One more word for getting around.
\end{grammarlens}

\subsection*{Your turn}
\begin{itemize}
\raggedright
  \item \emph{Say it:} \emph{tuìcháo}
  \item \emph{Say it:} \emph{tuìcháo}, once more
  \item \emph{Say it:} \emph{xiàwǔ tuìcháo de shíhou wǒmen qù hǎibiān zhǎo bèiké} — In the afternoon, when the tide goes out, we look for shells on the shore.
  \item \emph{Recall:} say the Mandarin for to go ashore, then the Mandarin for a seashell, then say \emph{tuìcháo} again
\end{itemize}

\begin{tcolorbox}[breakable,colback=teal!4,colframe=teal!35!black,title={Before you move on}]
What does tuìcháo mean? (\textquotedblleft{}To ebb (of the tide)\textquotedblright{}.) Say it once more.
\end{tcolorbox}
